\skripsiappendix{Contoh Materi Pendukung}
\label{lamp:contoh}

Lampiran ini memperlihatkan penempatan informasi pendukung di luar alur utama. Semua isinya merupakan petunjuk format, tanpa data penelitian pribadi. Pengguna dapat menggantinya dengan instrumen, rincian prosedur, daftar parameter, hasil pemeriksaan, atau bahan pendukung lain yang diperlukan pembaca dan diizinkan untuk dibagikan.

\section*{Keterangan Materi}

Setiap bahan pendukung perlu memiliki fungsi yang jelas. Jelaskan bagian naskah yang memerlukan bahan tersebut, asal bahan, serta batas penggunaannya. Apabila lampiran memuat salinan dokumen pihak lain, izin dan atribusi diperiksa sebelum salinan disertakan. Lampiran tidak sebaiknya menjadi tempat memasukkan seluruh berkas pengerjaan tanpa menilai kegunaannya bagi pembaca.

Tabel \ref{tab:contoh-lampiran} menunjukkan format catatan yang dapat diperiksa. Kolom tidak menyatakan bahwa suatu pemeriksaan sudah selesai. Pengguna dapat menambahkan status aktual serta tanggal hanya setelah ada bukti yang sesuai.

\begin{table}[htbp]
  \centering
  \caption{Contoh keterangan bahan pendukung}
  \label{tab:contoh-lampiran}
  \begin{tabular}{p{4.0cm}p{8.6cm}}
    \toprule
    Nama bahan & Keterangan yang disiapkan \\
    \midrule
    Sumber naskah & Nama berkas dan bagian yang menyimpan materi terkait \\
    Data bibliografis & Pengenal entri, sumber identitas, dan tanggal pemeriksaan \\
    Instrumen & Asal, tujuan, cara penggunaan, dan izin yang diperlukan \\
    Catatan perubahan & Bagian yang diperbaiki, alasan, dan dasar pemeriksaan \\
    Keluaran dokumen & Lingkungan penyusunan dan pemeriksaan tampilan terakhir \\
    \bottomrule
  \end{tabular}
\end{table}

\section*{Contoh Catatan Rujukan Silang}

Pola rujukan silang pada template dapat dipelajari melalui Tabel \ref{tab:contoh-keluaran}, Gambar \ref{fig:contoh-alur}, dan Persamaan \ref{eq:contoh-rerata}. Ketiganya mempunyai label unik yang mengikuti jenis unsur. Bab \ref{bab:hasil} juga merujuk lampiran ini, sehingga hubungan dari isi utama ke materi pendukung dan kembali ke isi utama dapat diperiksa.

Pengguna sebaiknya memberi nama label yang menggambarkan isi, lalu memanggil label ketika menyebut unsur dalam paragraf. Hindari memasukkan nomor tabel atau gambar secara manual karena nomor dapat berubah setelah penyuntingan. Nomor yang dibentuk otomatis tetap perlu diperiksa bersama makna kalimat yang menggunakannya.

\section*{Catatan sebelum Mengganti Contoh}

Simpan satu salinan sumber contoh untuk mempelajari pola penulisan. Pada naskah kerja, ganti identitas, abstrak, prakata, isi bab, dan bibliografi dengan materi sebenarnya. Hapus gambar atau tabel contoh apabila tidak lagi digunakan, lalu periksa rujukan yang sebelumnya menunjuk unsur tersebut. Materi yang dihapus tidak boleh meninggalkan nomor rujukan yang belum terdefinisi.

Penulis tetap bertanggung jawab terhadap isi ilmiah dan proses persetujuan naskah. Perangkat penyusunan membantu mengatur tampilan dan hubungan penanda, tetapi tidak memverifikasi kebenaran suatu temuan. Pemeriksaan akhir mencakup isi, sumber, tampilan, serta ketentuan akademik yang berlaku pada saat naskah diserahkan.
